
\documentclass[11pt,a4paper]{report}

% ==========================================
% CONFIGURACAO
% ==========================================

\usepackage{graphicx}


% =============================================
% CONFIGURACAO DO DOCUMENTO
% =============================================

\usepackage[utf8]{inputenc}
\usepackage[T1]{fontenc}
\usepackage[portuguese]{babel}
\usepackage{lmodern}
\usepackage{microtype}

\usepackage{geometry}
\geometry{
    a4paper,
    left=2.3cm,
    right=2.3cm,
    top=2.5cm,
    bottom=2.3cm
}

\usepackage[table]{xcolor}

% Cores
\definecolor{azulPDMI}{HTML}{1D5787}
\definecolor{azulClaro}{HTML}{EAF3FA}
\definecolor{cinzaLinha}{HTML}{D0D8DF}

\colorlet{azul}{azulPDMI}
\colorlet{azulclaro}{azulClaro}
\definecolor{verde}{HTML}{248563}

% Pacotes matematicos e graficos
\usepackage{amsmath,amssymb}
\usepackage{array,tabularx,longtable}
\usepackage{tikz}
\usetikzlibrary{arrows.meta}
\usepackage[most]{tcolorbox}

% Organizacao e formatacao
\usepackage{fancyhdr}
\usepackage{enumitem}
\usepackage{titlesec}
\usepackage{lastpage}
\usepackage{needspace}

% Espacamento geral
\setlength{\parindent}{0pt}
\setlength{\parskip}{4pt}
\setlength{\headheight}{26pt}
\renewcommand{\arraystretch}{1.55}

% =============================================
% CABECALHOS E RODAPES
% =============================================

\pagestyle{fancy}
\fancyhf{}

\newcommand{\currentparttitle}{}

\newcommand{\partwithheader}[1]{%
  \part{#1}%
  \gdef\currentparttitle{Parte \thepart\ -- #1}%
  \markboth{\currentparttitle}{}%
}

\renewcommand{\partmark}[1]{%
  \gdef\currentparttitle{Parte \thepart\ -- #1}%
  \markboth{\currentparttitle}{}%
}

\renewcommand{\chaptermark}[1]{%
  \markboth{\currentparttitle}
  {\chaptername\ \thechapter\ -- #1}%
}

\fancyhead[L]{%
  \small\textcolor{azulPDMI}{%
    \textbf{\nouppercase{\leftmark}}%
  }%
}

\fancyhead[R]{%
  \small\textcolor{azulPDMI}{%
    \textbf{\nouppercase{\rightmark}}%
  }%
}

\fancyfoot[R]{%
  \footnotesize Página \thepage\ de \pageref{LastPage}
}

\renewcommand{\headrulewidth}{0.4pt}
\renewcommand{\footrulewidth}{0pt}

\renewcommand{\headrule}{%
  \hbox to\headwidth{%
    \color{azulPDMI!35}%
    \leaders\hrule height \headrulewidth\hfill
  }%
}

\fancypagestyle{plain}{%
  \fancyhf{}
  \fancyhead[L]{%
    \small\textcolor{azulPDMI}{%
      \textbf{\nouppercase{\leftmark}}%
    }%
  }
  \fancyhead[R]{%
    \small\textcolor{azulPDMI}{%
      \textbf{\nouppercase{\rightmark}}%
    }%
  }
  \fancyfoot[R]{%
    \footnotesize Página \thepage\ de \pageref{LastPage}
  }
  \renewcommand{\headrulewidth}{0.4pt}
  \renewcommand{\footrulewidth}{0pt}
}

% =============================================
% FORMATACAO DOS CAPITULOS
% =============================================

\titleformat{\chapter}[display]
  {\normalfont\bfseries}
  {\large\chaptername\ \thechapter}
  {6pt}
  {\LARGE}

% Reduz o espaco acima e abaixo dos titulos
\titlespacing*{\chapter}
  {0pt}
  {3pt}
  {14pt}

% =============================================
% FORMATACAO DAS SECCOES
% =============================================

\titleformat{\section}
  {\large\bfseries\color{azulPDMI}}
  {}
  {0pt}
  {}

\titlespacing*{\section}
  {0pt}
  {10pt}
  {5pt}

% =============================================
% COMANDOS PERSONALIZADOS
% =============================================

\newcommand{\linhas}[1]{%
  \par\vspace{2pt}%
  \begingroup\color{cinzaLinha}%
  \foreach \n in {1,...,#1}{%
    \par\noindent\rule{\linewidth}{0.45pt}%
    \vspace{0.36cm}%
  }%
  \par\endgroup\vspace{2pt}%
}

\newcommand{\pergunta}[1]{%
  \textbf{#1}\enspace
}

\newcommand{\parte}[2]{%
  \section*{#1%
    \if\relax\detokenize{#2}\relax
    \else\ --- #2\fi
  }%
}

% =============================================
% CAIXAS
% =============================================

\newtcolorbox{desafio}{
  colback=azulClaro,
  colframe=azulPDMI!45,
  boxrule=0.5pt,
  arc=2mm,
  left=8pt,
  right=8pt,
  top=6pt,
  bottom=6pt
}

\newtcolorbox{nota}{
  colback=azulClaro,
  colframe=azulPDMI!35,
  boxrule=0.4pt,
  arc=1mm,
  left=8pt,
  right=8pt,
  top=6pt,
  bottom=6pt
}

\newtcolorbox{conclusao}{
  colback=azulClaro,
  colframe=azulPDMI!45,
  boxrule=0.5pt,
  arc=2mm,
  left=8pt,
  right=8pt,
  top=6pt,
  bottom=6pt
}

% =============================================
% POLIGONOS DA TABELA DA AULA 1
% =============================================

\newcommand{\poligonoTabela}[1]{%
  \begin{tikzpicture}[
    x=0.83cm,y=0.72cm,
    baseline=(current bounding box.center),
    line join=round
  ]
    \ifcase#1\relax
    \or\or\or
      \draw[draw=azulPDMI,fill=azulClaro,thick]
        (0.15,0.15)--(3.30,0.12)--(1.65,2.10)--cycle;
    \or
      \draw[draw=azulPDMI,fill=azulClaro,thick]
        (0.15,0.20)--(3.30,0.15)--(3.20,2.15)--(0.40,2.30)--cycle;
    \or
      \draw[draw=azulPDMI,fill=azulClaro,thick]
        (0.15,0.55)--(1.20,0.08)--(3.10,0.40)--
        (3.15,1.85)--(1.45,2.35)--cycle;
    \or
      \draw[draw=azulPDMI,fill=azulClaro,thick]
        (0.30,0.55)--(1.10,0.07)--(2.55,0.15)--
        (3.35,1.10)--(2.75,2.25)--(0.65,2.15)--cycle;
    \or
      \draw[draw=azulPDMI,fill=azulClaro,thick]
        (0.15,0.85)--(0.55,0.12)--(1.85,0.08)--
        (3.20,0.75)--(3.05,1.75)--(2.05,2.35)--
        (0.55,2.20)--cycle;
    \fi
  \end{tikzpicture}%
}

% Designacao dos capitulos
\addto\captionsportuguese{%
  \renewcommand{\chaptername}{Secção}%
}


\begin{document}

% ==========================================
% CAPA
% ==========================================

\begin{titlepage}

\thispagestyle{empty}

\begin{tikzpicture}[remember picture,overlay]

% Faixa azul superior

\fill[azulPDMI]
(current page.north west)
rectangle
([yshift=-1.1cm]current page.north east);

% Linha decorativa inferior

\draw[azulPDMI!30,line width=0.8pt]
([xshift=2.5cm,yshift=1.6cm]current page.south west)
--
([xshift=-2.5cm,yshift=1.6cm]current page.south east);

% ==========================================
% LOGOTIPO DA NOVA FCT
% ==========================================


% ==========================================
% LOGOTIPO NOVA FCT
% ==========================================

\node[
    anchor=north west,
    xshift=2.3cm,
    yshift=-1.5cm,
    inner sep=0pt
] at (current page.north west)
{
    \includegraphics[
        width=3.2cm,
        keepaspectratio
    ]{LOGO.png}
};


\end{tikzpicture}

\vspace*{2.4cm}

\begin{center}

% ==========================================
% IDENTIFICACAO INSTITUCIONAL
% ==========================================

{\large\bfseries\color{azulPDMI}
UNIVERSIDADE NOVA DE LISBOA}

\par\vspace{0.15cm}

{\Large\bfseries NOVA FCT}

\par\vspace{0.08cm}

{\normalsize
Faculdade de Ciências e Tecnologia}

\par\vspace{0.55cm}

{\normalsize
Mestrado em Ensino de Matemática}

\par\vspace{0.12cm}

{\normalsize\bfseries
Pedagogia e Didática da Matemática I}

% ==========================================
% TITULO DO TRABALHO
% ==========================================

\par\vspace{1.25cm}

{\Huge\bfseries\color{azulPDMI}
Trabalho Prático}

\par\vspace{0.35cm}

{\Large
Dois caminhos para a mesma propriedade}

\par\vspace{0.4cm}

{\large\color{black!70}
Soma das amplitudes dos ângulos internos\\
de polígonos convexos}

\par\vspace{0.2cm}

{\normalsize
Matemática -- 7.º ano}

% ==========================================
% REPRESENTACOES GEOMETRICAS
% ==========================================

\par\vspace{0.85cm}

\begin{tikzpicture}[
    line join=round,
    scale=0.85,
    >=Stealth
]

% PRIMEIRA ABORDAGEM
% Decomposicao em triangulos

\coordinate (A) at (-5.2,0.1);
\coordinate (B) at (-2.8,0.1);
\coordinate (C) at (-2.15,1.6);
\coordinate (D) at (-3.55,2.6);
\coordinate (E) at (-5.6,1.55);

\draw[
    fill=azulClaro,
    draw=azulPDMI,
    line width=1.2pt
]
(A)--(B)--(C)--(D)--(E)--cycle;

% Diagonais

\draw[
    azulPDMI!70,
    densely dashed,
    line width=1.1pt
]
(A)--(C)
(A)--(D);

% Vertices

\foreach \p in {A,B,C,D,E}
    \fill[azulPDMI] (\p) circle (1.8pt);

\node[
    font=\footnotesize,
    text=azulPDMI,
    align=center
]
at (-3.83,-0.55)
{Decomposição em triângulos};

% ==========================================
% SEGUNDA ABORDAGEM
% Mudancas de direcao
% ==========================================

\coordinate (F) at (0.95,0.15);
\coordinate (G) at (4.05,0.15);
\coordinate (H) at (4.65,2.37);
\coordinate (I) at (1.58,2.37);

\draw[
    fill=azulClaro,
    draw=azulPDMI,
    line width=1.2pt
]
(F)--(G)--(H)--(I)--cycle;

% Setas do percurso

\draw[
    -{Stealth[length=2.8mm]},
    azulPDMI,
    line width=1.7pt
]
(1.55,0.15)--(3.26,0.15);

\draw[
    -{Stealth[length=2.8mm]},
    azulPDMI,
    line width=1.7pt
]
(4.23,0.82)--(4.47,1.77);

\draw[
    -{Stealth[length=2.8mm]},
    azulPDMI,
    line width=1.7pt
]
(4.03,2.37)--(2.30,2.37);

\node[
    font=\footnotesize,
    text=azulPDMI,
    align=center
]
at (2.84,-0.55)
{Mudanças de direção};

\end{tikzpicture}

% ==========================================
% IDENTIFICACAO DAS DUAS AULAS
% ==========================================

\par\vspace{0.9cm}

\colorbox{azulClaro}{%
    \parbox{0.80\textwidth}{%
        \centering
        \vspace{0.2cm}

        {\bfseries\color{azulPDMI}
        Aula simulada e aula alternativa}

        \par\vspace{0.1cm}

        Duas abordagens didáticas ao mesmo
        conteúdo matemático

        \vspace{0.2cm}
    }%
}

\vfill

% ==========================================
% ELEMENTOS DO GRUPO
% ==========================================

{\bfseries
Elementos do grupo}

\par\vspace{0.2cm}

\begin{tabular}{@{}l@{\qquad}l@{}}

Maria Botelho
&
N.º 79842
\\[6pt]

Rui Costa
&
N.º 79463
\\[6pt]

Luiz Júnior
&
N.º
\end{tabular}

\par\vspace{0.55cm}

{\small\color{black!65}
Ano letivo 2026/2027}

\end{center}

\end{titlepage}

% ==========================================
% INDICE
% ==========================================

\setcounter{secnumdepth}{3}
\setcounter{tocdepth}{3}

\begingroup

\renewcommand{\currentparttitle}{}

\makeatletter
\let\@mkboth\@gobbletwo
\makeatother

\tableofcontents

\endgroup

\clearpage

% ==========================================
% PARTE I - AULA SIMULADA
% ==========================================

\partwithheader{Aula simulada}

% PLANIFICACAO

\chapter{Planificação}


% Planificacao da Aula 1

\begin{center}
{\large\bfseries\color{azulPDMI}Quantos triângulos se escondem num polígono?}\\[3pt]
{\small Planificação da aula - Matemática, 7.º ano}
\end{center}

\section*{1. Objeto}
\textbf{Tema:} Geometria. \textbf{Tópico:} Soma das amplitudes dos ângulos internos de um polígono convexo.\\
\textbf{Objeto matemático central:} dedução e justificação da soma das amplitudes dos ângulos internos de um polígono convexo com $n$ lados.\\
\textbf{Proposta de duração:} 50 minutos.

\section*{2. Objetivos}
\subsection*{Objetivos gerais}
A presente aula deve ajudar a desenvolver as capacidades de:
\begin{itemize}[leftmargin=*,itemsep=3pt]
\item \textbf{Raciocínio matemático.}
\item \textbf{Representações matemáticas:} usar a representação algébrica e a representação gráfica e fazer conexões entre estas representações.
\item \textbf{Trabalho colaborativo:} trabalhar em grupo.
\item \textbf{Comunicação matemática:} comunicar conteúdos matemáticos.
\end{itemize}

\subsection*{Objetivos específicos}
No final da aula, o aluno deve ser capaz de:
\begin{enumerate}[label=\textbf{OE\arabic*.},leftmargin=*,itemsep=4pt]
\item Decompor um polígono convexo em triângulos não sobrepostos, através do desenho de diagonais que não se intersetam.
\item Reconhecer a relação algébrica entre o número de lados ($n$) de um polígono convexo e o número de triângulos ($m$) em que se decompõe $\left[\,m=n-2\,\right]$.
\item Reconhecer a relação algébrica entre o número de lados ($n$) de um polígono convexo e a soma ($S$) das amplitudes dos seus ângulos internos $\left[\,S=(n-2)\times180^\circ\,\right]$.
\end{enumerate}

\section*{3. Pré-requisitos}
Ao iniciar a aula, o aluno deverá:
\begin{itemize}[leftmargin=*,itemsep=2pt]
\item Ter presente os conceitos de:
\begin{itemize}[itemsep=1pt]
\item vértices consecutivos ou adjacentes;
\item polígono convexo.
\end{itemize}
\item Ter presente que a soma das amplitudes dos ângulos internos de um triângulo é $180^\circ$.
\item Ser capaz de resolver equações simples.
\end{itemize}

\section*{4. Recursos e materiais}
\begin{itemize}[leftmargin=*,itemsep=2pt]
\item \textbf{A fornecer pelo professor:} folha em papel com o enunciado, a distribuir a cada grupo de trabalho.
\item \textbf{A trazer pelo aluno:}  lápis e borracha.
\item \textbf{Opcional:} régua e papel de rascunho.
\end{itemize}

\section*{5. Desenvolvimento da aula (50 minutos)}
\begingroup\small
\renewcommand{\arraystretch}{1.24}
\begin{longtable}{|p{1.35cm}|p{2.25cm}|p{5.55cm}|p{5.1cm}|}
\hline
\rowcolor{azulClaro}\textbf{Tempo}&\textbf{Fase / questões}&\textbf{Intervenção do professor}&\textbf{Atividade dos alunos}\\\hline
\endfirsthead
\hline\rowcolor{azulClaro}\textbf{Tempo}&\textbf{Fase / questões}&\textbf{Intervenção intencional do professor}&\textbf{Atividade dos alunos}\\\hline
\endhead
\textbf{0--5 min} (5 min)&Introdução. &
Apresenta a tarefa (tema e tópico). Recorda que a soma dos ângulos internos de um triângulo é 180º. Relembra o conceito de figura geométrica convexa. Coordena a organização a pares. &Leem o desafio, recordam que a soma dos ângulos internos de um triângulo é $180^\circ$ e interpretam «sem medir nenhum ângulo».\\\hline
\textbf{5--13 min} (8 min)&Parte I; 1.1&Circula e escuta antes de intervir. Regista as estratégias (tentativas, recorte de ângulos, decomposição em triângulos). Pergunta «O que já sabes que te possa ajudar?» sem sugerir a diagonal.&Procuram, a pares, uma estratégia para obter a soma dos ângulos internos do quadrilátero e do pentágono sem transferidor; desenham sobre as figuras e registam o raciocínio.\\\hline
\textbf{13--23 min} (10 min)&Parte I; 1.2, caixa de estratégia e 1.3&Introduz «diagonal» só depois da comparação. Pergunta «Há sobreposições ou regiões por preencher?». Apoia quem não relaciona os ângulos dos triângulos com os do polígono pedindo que os assinale com cores. A 1.3 fica essencialmente oral.&Comparam estratégias com o colega; traçam diagonais no quadrilátero e no pentágono; calculam $2\times180^\circ$ e $3\times180^\circ$; explicam, oralmente ou em registo curto, porque se podem somar os ângulos dos triângulos.\\\hline
\textbf{23--32 min} (9 min)&Parte II; 2.1--2.3&Insiste em traçar as diagonais a partir de um só vértice. Se surgirem diagonais cruzadas, pergunta «Os triângulos ficam sem sobreposições?». Pede: «O que representa $n-2$?» «Porquê $180^\circ$?». Não corrige de imediato o erro de $n$ triângulos.&Completam a tabela (quadrilátero e pentágono já estudados; hexágono novo), descrevem a relação entre lados e triângulos e propõem $S=(n-2)\times180^\circ$, explicando os fatores.\\\hline
\textbf{32--38 min} (6 min)&Parte II; 2.4&Pergunta «Quantos vértices não podes ligar ao vértice escolhido?» (ele próprio e os dois adjacentes). Distingue verificar na tabela de justificar para $n$ qualquer.&Justificam que a regra vale para qualquer polígono convexo, pensando nas $n-3$ diagonais a partir de um vértice e nos $n-2$ triângulos.\\\hline
\textbf{38--46 min} (8 min)&Discussão coletiva&Sequencia: estratégia por tentativa; decomposição do quadrilátero; tabela e regularidade; justificação com $n-3$ e $n-2$. Pede que expliquem cada fator antes de simplificar. «E se não fosse convexo?» fica como ponto de reflexão, sem aprofundar.&Apresentam 2 ou 3 resoluções previamente selecionadas; comparam estratégias; discutem o significado de $n-2$ e porque a decomposição funciona num polígono convexo.\\\hline
\textbf{46--50 min} (4 min)&Sistematização e conclusão&Institucionaliza a regra e a condição de convexidade. Recolhe as respostas individuais e identifica o que retomar na aula seguinte.&Completam individualmente «Para concluir» e determinam a soma para 12 lados ($1800^\circ$), justificando.\\\hline
\end{longtable}
\endgroup

\section*{6. Antecipação das dificuldades e respostas do professor}
\begingroup\small
\renewcommand{\arraystretch}{1.22}
\begin{longtable}{|p{2.55cm}|p{5.7cm}|p{6.65cm}|}
\hline
\rowcolor{azulClaro}\textbf{Questão / obstáculo}&\textbf{Produção ou dificuldade previsível}&\textbf{Intervenção docente planeada}\\\hline
\endfirsthead
\hline\rowcolor{azulClaro}\textbf{Questão / obstáculo}&\textbf{Produção ou dificuldade previsível}&\textbf{Intervenção docente planeada}\\\hline
\endhead
1.1: estratégia inicial&Medir os ângulos com o transferidor, apesar do pedido; recortar e juntar os ângulos; ou dizer que «não se consegue sem medir». Há ainda quem proponha uma regra sem fundamento (por exemplo, $360^\circ$ para qualquer figura).&«O que já sabes sobre a soma dos ângulos de um triângulo?» «Consegues ver triângulos dentro destas figuras?» Valorizar o recorte de ângulos como ideia, sem validar $360^\circ$ para todos. Não sugerir ainda a diagonal.\\\hline
1.1: desenho sobre as figuras&Traçar segmentos que não unem vértices, ou desenhar triângulos que não ocupam toda a figura, deixando regiões por preencher.&«Os extremos dos teus segmentos são vértices?» «Toda a região do polígono ficou coberta?» Pedir que marquem os triângulos com cores diferentes.\\\hline
1.2: comparação de estratégias&Os pares trocam resultados sem justificar; ou consideram que estratégias diferentes que dão o mesmo valor são equivalentes sem as relacionar.&«Porque é que as duas estratégias dão o mesmo resultado?» Pedir que cada elemento do par explique a ideia do colega. Registar as estratégias para a discussão.\\\hline
Caixa de estratégia, a)--b)&Traçar a diagonal do quadrilátero e contar 2 triângulos; no pentágono, traçar diagonais de vértices diferentes, que se cruzam, e obter mais de 3 triângulos; ou confundir diagonal com lado.&Recordar a definição: «Une vértices consecutivos?» Pedir o regresso a um único vértice e a contagem de triângulos sem sobreposições. Mostrar, se necessário, um exemplo com diagonais cruzadas e perguntar se os triângulos se sobrepõem.\\\hline
Caixa de estratégia, c)&Calcular $2\times180^\circ=360^\circ$ e $3\times180^\circ=540^\circ$ e não perceber que se trata da soma dos ângulos internos do polígono; ou somar $180^\circ$ por cada vértice (obter $720^\circ$ no quadrilátero).&«Que ângulos estão dentro de cada triângulo?» «Eles coincidem com os do quadrilátero?» Pedir que marquem, com a mesma cor, os ângulos do triângulo e a parte correspondente do ângulo do polígono.\\\hline
1.3: justificação da soma&Responder «porque dá certo» ou «porque o professor disse»; ou não reconhecer que, num vértice, o ângulo do polígono se divide em ângulos de dois triângulos.&«Há ângulos de triângulos que ficaram de fora ou que se repetem?» «O que acontece no vértice de onde partem as diagonais?» Aceitar respostas orais e registos curtos, e exigir referência às sobreposições e regiões por preencher.\\\hline
2.1: preenchimento da tabela&Traçar diagonais de vértices diferentes no hexágono; contar $n$ triângulos (6 no hexágono) em vez de $n-2$; esquecer-se de calcular a soma ou escrever $4\times180^\circ=620^\circ$ por erro de cálculo.&«De onde partem todas as diagonais?» «Quantos vértices não consegues ligar ao vértice escolhido?» Confrontar com o desenho: contar os triângulos um a um. Não indicar o valor; pedir que verifiquem a multiplicação.\\\hline
2.2: relação lados--triângulos&Dizer apenas «aumenta um de cada vez»; ou descrever a relação de forma recursiva («soma-se mais um triângulo») sem relacionar $n$ com o número de triângulos.&«Se tiver 10 lados, quantos triângulos terá?» Pedir que a descrição permita responder sem desenhar. Valorizar a ideia de «menos dois» e pedir que a justifiquem pelo desenho.\\\hline
2.3: simbolização&Escrever $S=n\times180^\circ$; escrever $(n-2)\times180^\circ$ por padrão, sem interpretar os fatores; ou escrever $180^\circ-2n$ ou $n-2\times180^\circ$ (sem parênteses).&«O que representa $n-2$?» «O que representa $180^\circ$?» Testar a expressão no quadrilátero e no pentágono. Se faltarem parênteses, pedir que calculem com $n=5$ e comparem com o valor da tabela.\\\hline
2.4: generalização&«Funciona no triângulo, quadrilátero, pentágono e hexágono, portanto funciona sempre»; ou explicar apenas com os exemplos da tabela, sem argumentar para $n$ qualquer.&«Quantos exemplos estudámos?» «Que argumento não depende do número de exemplos?» «Quantos vértices não podes ligar ao vértice escolhido?» (ele próprio e os dois adjacentes, logo $n-3$ diagonais e $n-2$ triângulos).\\\hline
2.4: convexidade&Não perceber porque o polígono tem de ser convexo; ou pensar que a regra só vale para polígonos regulares.&«Num polígono não convexo, as diagonais ficam sempre dentro do polígono?» Desenhar um exemplo não convexo como ponto de reflexão, sem aprofundar. Esclarecer que a regra não exige lados ou ângulos iguais.\\\hline
Conclusão e 12 lados&Copiar a expressão sem completar o «porque»; ou calcular $12\times180^\circ=2160^\circ$, esquecendo-se de subtrair 2; ou desenhar o dodecágono em vez de aplicar a regra.&«Quantos triângulos terias com 12 lados?» «Onde está esse número na tua expressão?» Pedir que expliquem o «porque» em linguagem própria. Recolher as respostas individuais para identificar o que retomar.\\\hline
Dificuldade técnica e de gestão&Material desenhado em escala pequena, dificultando o traçado; poucos alunos a terminar a tabela; discussões longas no trabalho a pares que reduzem o tempo da discussão coletiva.&Ter cópias adicionais das figuras e cores disponíveis; admitir que 1.3 e 2.4 sejam orais; avisar o tempo restante; preservar os últimos 4 minutos para o registo individual.\\\hline
\end{longtable}
\endgroup

\section*{7. Avaliação das aprendizagens}
\textbf{Instrumentos:} observação durante o trabalho a pares, respostas escritas no enunciado (incluindo «Para concluir»), argumentação na discussão coletiva e resposta individual ao caso dos 12 lados.

\textbf{Critério de resposta completa:} o aluno explicita que, a partir de um vértice, se traçam $n-3$ diagonais, que ficam $n-2$ triângulos, cada um com soma de ângulos internos $180^\circ$, pelo que $S=(n-2)\times180^\circ$; e aplica-o a 12 lados, obtendo $1800^\circ$. Respostas que apresentem apenas a expressão correta, sem justificação, revelam aprendizagem procedimental ainda incompleta.

% ENUNCIADO

\chapter{Enunciado}


\begin{center}
{\LARGE\bfseries\color{azulPDMI}
Quantos triângulos se escondem num polígono?}\\[2pt]
{\small Soma das amplitudes dos ângulos internos
-- tarefa de exploração}
\end{center}

\vspace{2pt}

\noindent
\textbf{Nome:} \rule{8.1cm}{0.4pt}
\hfill
\textbf{N.º:} \rule{1.3cm}{0.4pt}

\begin{tcolorbox}[
colback=azulClaro,
colframe=azulPDMI!45,
boxrule=0.5pt,
arc=2mm,
left=7pt,
right=7pt,
top=5pt,
bottom=5pt
]

\textbf{O desafio}\par

Sabemos que a soma das amplitudes dos
ângulos internos de um triângulo é
$180^\circ$.

Conseguiremos descobrir a soma dos
ângulos internos de um
\textbf{polígono convexo sem medir
nenhum ângulo}?

Regista desenhos, cálculos e explicações
que permitam compreender o teu raciocínio.

\end{tcolorbox}

% ========================================================
% PARTE I
% ========================================================

\section*{Parte I -- À procura de uma estratégia}

\textbf{1.1.}
Observa o quadrilátero e o pentágono.
Procura uma estratégia para obter a soma
das amplitudes dos ângulos internos de
cada um, \textbf{sem utilizar o transferidor}.

Podes desenhar sobre as figuras.
Explica como pensaste.

\vspace{2pt}

\begin{center}

\begin{minipage}{.45\linewidth}
\centering

\begin{tikzpicture}[scale=.65,line join=round]

\coordinate (A) at (.3,.2);
\coordinate (B) at (4.3,.1);
\coordinate (C) at (4.7,2.6);
\coordinate (D) at (.5,2.9);

\draw[
thick,
draw=azulPDMI,
fill=azulClaro
]
(A)--(B)--(C)--(D)--cycle;

\node[below left] at (A) {$A$};
\node[below right] at (B) {$B$};
\node[above right] at (C) {$C$};
\node[above left] at (D) {$D$};

\end{tikzpicture}

\vspace{2pt}

{\small Quadrilátero $ABCD$}

\end{minipage}
\hfill
\begin{minipage}{.45\linewidth}
\centering

\begin{tikzpicture}[scale=.60,line join=round]

\coordinate (A) at (.4,.2);
\coordinate (B) at (3.8,.1);
\coordinate (C) at (5,2.2);
\coordinate (D) at (2.8,3.4);
\coordinate (E) at (.1,2);

\draw[
thick,
draw=azulPDMI,
fill=azulClaro
]
(A)--(B)--(C)--(D)--(E)--cycle;

\node[below left] at (A) {$A$};
\node[below right] at (B) {$B$};
\node[right] at (C) {$C$};
\node[above] at (D) {$D$};
\node[left] at (E) {$E$};

\end{tikzpicture}

\vspace{2pt}

{\small Pentágono $ABCDE$}

\end{minipage}

\end{center}

\vspace{-5pt}

\linhas{2}

\textbf{1.2.}
Compara a tua ideia com a de um colega.
Chegaram ao mesmo resultado?
Identificaram estratégias diferentes?
Partilhem uma explicação.

\linhas{1}

\begin{tcolorbox}[
colback=azulClaro,
colframe=azulPDMI!45,
boxrule=.5pt,
arc=2mm,
left=7pt,
right=7pt,
top=5pt,
bottom=5pt
]

\textbf{Uma estratégia para experimentar}

Uma \textbf{diagonal} une dois vértices
não consecutivos de um polígono.

Mesmo que tenhas encontrado outra
estratégia, experimenta agora:

\begin{enumerate}[
label=\alph*),
leftmargin=1.7em,
nosep,
topsep=2pt
]

\item
Traça uma diagonal no quadrilátero.
Em que figuras ficou dividido?

\item
No pentágono, escolhe um vértice
e traça todas as diagonais que
partem dele.
Quantos triângulos obténs?

\item
Usa a soma dos ângulos internos
de um triângulo para calcular a
soma dos ângulos internos de
cada polígono.

\end{enumerate}

\end{tcolorbox}

\textbf{1.3.}
Porque é que podemos somar os ângulos
dos triângulos para obter a soma dos
ângulos internos do polígono?

Na tua explicação, considera se há
sobreposições ou regiões por preencher.

\linhas{2}


\section*{Parte II -- Da exploração à regra geral}

\textbf{2.1.}
Completa a tabela.

Em cada polígono, desenha as diagonais
\textbf{a partir de um único vértice},
de modo a obter triângulos sem
sobreposições e sem espaços por preencher.

\begin{center}

\renewcommand{\arraystretch}{1.12}

\begin{tabularx}{\linewidth}{|
>{\centering\arraybackslash}p{3.5cm}|
>{\centering\arraybackslash}p{1.6cm}|
>{\centering\arraybackslash}p{2.2cm}|
>{\centering\arraybackslash}X|}

\hline

\rowcolor{azulClaro}

\textbf{Polígono} &
\textbf{N.º de lados} &
\textbf{N.º de triângulos} &
\textbf{Soma dos ângulos internos}

\\ \hline

% TRIANGULO

\parbox[c][1.92cm][c]{3.5cm}{
\centering

\begin{tikzpicture}[
scale=.72,
baseline=(current bounding box.center)
]

\draw[
thick,
azulPDMI,
fill=azulClaro
]
(0,0)--(2.8,0)--(1.4,1.7)--cycle;

\end{tikzpicture}

\vspace{3pt}

{\scriptsize Triângulo}
}

& 3 & 1 & $180^\circ$

\\ \hline

% QUADRILATERO

\parbox[c][1.92cm][c]{3.5cm}{
\centering

\begin{tikzpicture}[
scale=.69,
baseline=(current bounding box.center)
]

\draw[
thick,
azulPDMI,
fill=azulClaro
]
(0,0)--(2.8,.1)--(2.7,1.8)--
(.2,1.7)--cycle;

\end{tikzpicture}

\vspace{3pt}

{\scriptsize Quadrilátero}
}

& 4 & &

\\ \hline

% PENTAGONO

\parbox[c][1.92cm][c]{3.5cm}{
\centering

\begin{tikzpicture}[
scale=.69,
baseline=(current bounding box.center)
]

\draw[
thick,
azulPDMI,
fill=azulClaro
]
(0,.6)--(1,.05)--(2.7,.2)--
(3,1.6)--(1.2,2)--cycle;

\end{tikzpicture}

\vspace{3pt}

{\scriptsize Pentágono}
}

& 5 & &

\\ \hline

% HEXAGONO

\parbox[c][1.92cm][c]{3.5cm}{
\centering

\begin{tikzpicture}[
scale=.69,
baseline=(current bounding box.center)
]

\draw[
thick,
azulPDMI,
fill=azulClaro
]
(.2,.4)--(1,.05)--(2.4,.2)--
(3,1.1)--(2.1,2)--(.3,1.7)--cycle;

\end{tikzpicture}

\vspace{3pt}

{\scriptsize Hexágono}
}

& 6 & &

\\ \hline

\end{tabularx}

\end{center}

\renewcommand{\arraystretch}{1.55}

\textbf{2.2.}
Que relação observas entre o número
de lados do polígono e o número
de triângulos obtidos?

Descreve-a por palavras.

\linhas{1}

\textbf{2.3.}
Considera um polígono convexo
com $n$ lados.

Propõe uma expressão para a soma
$S$ das amplitudes dos seus
ângulos internos.

Explica o significado dos
fatores utilizados.

\linhas{2}
\newpage
\textbf{2.4.}
A tua regra funciona para
\emph{qualquer} polígono convexo?

Procura justificá-la sem te baseares
apenas nos exemplos da tabela.

Pensa no número de diagonais que
podes traçar a partir de um vértice
e no número de triângulos assim formado.

\linhas{2}

% ========================================================
% CONCLUSAO
% ========================================================

\begin{tcolorbox}[
colback=azulClaro,
colframe=azulPDMI!40,
boxrule=.5pt,
arc=2mm,
left=7pt,
right=7pt,
top=5pt,
bottom=5pt
]

\textbf{Para concluir}

Completa:

«Num polígono convexo com $n$ lados,
a soma das amplitudes dos ângulos
internos é

\[
S=\underline{\hspace{3.0cm}}
\]

porque
\underline{\hspace{4cm}}».

\end{tcolorbox}

\vspace{2pt}

{\footnotesize\itshape

\textbf{Se houver tempo:}

Sem desenhar, determina a soma das
amplitudes dos ângulos internos
de um polígono convexo com 12 lados.

Justifica.

}


% RESOLUCAO

\chapter{Resolução}

\begingroup
\small
\setlength{\parskip}{2pt}
\setlength{\abovedisplayskip}{5pt}
\setlength{\belowdisplayskip}{5pt}
\begin{center}
{\Large\bfseries\color{azulPDMI}
Quantos triângulos se escondem num polígono?}\\[3pt]
{\small Resolução da tarefa de exploração -- aula simulada (50 minutos)}
\end{center}

\begin{tcolorbox}[
colback=azulClaro,colframe=azulPDMI!40,boxrule=.5pt,
arc=2mm,left=8pt,right=8pt,top=6pt,bottom=6pt]
\textbf{Ideia matemática central.}
A propriedade da soma dos ângulos internos de um triângulo
($180^\circ$) permite determinar a soma dos ângulos internos de um
polígono convexo, desde que este seja decomposto em triângulos que
preencham toda a sua região interior sem sobreposições.
A exploração dos exemplos conduz à expressão
\[
\boxed{S=(n-2)\times180^\circ,\qquad n\geq3,}
\]
cuja validade geral exige uma justificação matemática.
\end{tcolorbox}

\section*{Parte I -- À procura de uma estratégia}

\subsection*{1.1. Quadrilátero e pentágono}

Uma estratégia possível é traçar, a partir de um vértice, diagonais
que decomponham o polígono em triângulos. Os desenhos seguintes
ilustram essa estratégia.

\begin{center}
\begin{minipage}{.47\linewidth}
\centering
\begin{tikzpicture}[scale=.67,line join=round]
\coordinate (A) at (.3,.2);
\coordinate (B) at (4.3,.1);
\coordinate (C) at (4.7,2.6);
\coordinate (D) at (.5,2.9);
\draw[thick,draw=azulPDMI,fill=azulClaro]
  (A)--(B)--(C)--(D)--cycle;
\draw[very thick,orange!85!black] (A)--(C);
\node[below left] at (A) {$A$};
\node[below right] at (B) {$B$};
\node[above right] at (C) {$C$};
\node[above left] at (D) {$D$};
\end{tikzpicture}\\[3pt]
{\small Quadrilátero: 2 triângulos}
\end{minipage}\hfill
\begin{minipage}{.47\linewidth}
\centering
\begin{tikzpicture}[scale=.64,line join=round]
\coordinate (A) at (.4,.2);
\coordinate (B) at (3.8,.1);
\coordinate (C) at (5,2.2);
\coordinate (D) at (2.8,3.4);
\coordinate (E) at (.1,2);
\draw[thick,draw=azulPDMI,fill=azulClaro]
  (A)--(B)--(C)--(D)--(E)--cycle;
\draw[very thick,orange!85!black] (A)--(C) (A)--(D);
\node[below left] at (A) {$A$};
\node[below right] at (B) {$B$};
\node[right] at (C) {$C$};
\node[above] at (D) {$D$};
\node[left] at (E) {$E$};
\end{tikzpicture}\\[3pt]
{\small Pentágono: 3 triângulos}
\end{minipage}
\end{center}

No quadrilátero $ABCD$, a diagonal $AC$ forma os triângulos
$ABC$ e $ACD$. Portanto,
\[
S_4=2\times180^\circ=\boxed{360^\circ}.
\]

No pentágono $ABCDE$, as diagonais $AC$ e $AD$ formam os
triângulos $ABC$, $ACD$ e $ADE$. Assim,
\[
S_5=3\times180^\circ=\boxed{540^\circ}.
\]

\textbf{Nota:} Esta é uma resolução possível, não a única.
Podem escolher-se outros vértices e outras diagonais, desde que se
obtenha uma decomposição válida. Um procedimento diferente, como
ligar um ponto interior aos vértices, exige ter em conta os ângulos
adicionais em torno desse ponto.

\subsection*{1.2. Comparação com a estratégia de um colega}

\textbf{Exemplo de resposta possível:}
«Escolhemos diagonais diferentes, mas obtivemos os mesmos resultados:
$360^\circ$ no quadrilátero e $540^\circ$ no pentágono.
Em ambos os casos, dividimos a figura em triângulos e utilizámos a
soma de $180^\circ$ em cada triângulo.»

A resposta é aberta: os alunos podem ter utilizado estratégias distintas
ou identificado dúvidas diferentes. Importa comparar os procedimentos,
explicar por que são válidos e não apenas verificar se os resultados
numéricos coincidem.

\Needspace{9\baselineskip}
\subsection*{Uma estratégia para experimentar -- alíneas a), b) e c)}

\textbf{a)} A diagonal $AC$ do quadrilátero origina dois triângulos,
$ABC$ e $ACD$.

\textbf{b)} No pentágono, escolhendo $A$ como vértice de partida,
traçam-se as diagonais $AC$ e $AD$, obtendo-se três triângulos:
$ABC$, $ACD$ e $ADE$.

\textbf{c)} Cada triângulo tem soma das amplitudes dos ângulos
internos igual a $180^\circ$. Logo,
\[
\begin{aligned}
S_4&=2\times180^\circ=360^\circ,\\
S_5&=3\times180^\circ=540^\circ.
\end{aligned}
\]

\subsection*{1.3. Justificação da soma dos ângulos dos triângulos}

As diagonais traçadas no interior dos polígonos convexos
originam triângulos cujas regiões interiores não se sobrepõem
e cuja reunião é exatamente a região interior de cada polígono.

Além disso, \textbf{os ângulos dos triângulos que têm o mesmo vértice
do polígono recompõem, sem lacunas nem sobreposições, o ângulo
interno desse vértice}. Por exemplo, no quadrilátero,
\[
\widehat{DAB}=\widehat{DAC}+\widehat{CAB}
\quad\text{e}\quad
\widehat{BCD}=\widehat{BCA}+\widehat{ACD}.
\]
Nos vértices $B$ e $D$, os ângulos dos triângulos coincidem
com os ângulos internos correspondentes do quadrilátero.

Assim, a soma de todos os ângulos internos dos triângulos é
\textbf{exatamente} a soma dos ângulos internos do polígono;
não são acrescentados ângulos novos em pontos interiores, porque
os segmentos utilizados têm extremidades nos vértices do polígono.

\section*{Parte II -- Da exploração à regra geral}

\subsection*{2.1. Decomposição e preenchimento da tabela}

Traçando as diagonais a partir de um mesmo vértice, obtém-se:

\begin{center}
\renewcommand{\arraystretch}{1.24}
\begin{tabularx}{\linewidth}{|
>{\raggedright\arraybackslash}X|
>{\centering\arraybackslash}p{2.0cm}|
>{\centering\arraybackslash}p{2.5cm}|
>{\centering\arraybackslash}p{3.25cm}|}
\hline
\rowcolor{azulClaro}
\textbf{Polígono} &
\textbf{N.º de lados} &
\textbf{N.º de triângulos} &
\textbf{Soma dos ângulos internos} \\\hline
Triângulo & 3 & 1 & $1\times180^\circ=180^\circ$ \\\hline
Quadrilátero & 4 & 2 & $2\times180^\circ=360^\circ$ \\\hline
Pentágono & 5 & 3 & $3\times180^\circ=540^\circ$ \\\hline
Hexágono & 6 & 4 & $4\times180^\circ=720^\circ$ \\\hline
\end{tabularx}
\end{center}

\textbf{Construções esperadas:} no triângulo não é preciso
traçar diagonais; no quadrilátero traça-se uma diagonal;
no pentágono, duas; no hexágono, três. Em todos os casos,
as diagonais partem de um único vértice e produzem,
respetivamente, $1$, $2$, $3$ e $4$ triângulos.

\subsection*{2.2. Relação entre o número de lados e o número de triângulos}

\textbf{Resposta possível:}
«O número de triângulos obtidos é sempre inferior em duas unidades
ao número de lados do polígono. Quando passamos de um polígono
de $n$ lados para outro de $n+1$ lados, obtemos mais um triângulo
pela mesma estratégia de decomposição.»

Se $m$ designar o número de triângulos, a relação é
\[
\boxed{m=n-2.}
\]
Nesta fase, os exemplos permitem \emph{conjeturar} a relação.
A questão 2.4 destina-se a justificar que ela é válida em geral.

\subsection*{2.3. Expressão para a soma dos ângulos internos}

Se um polígono convexo com $n$ lados se decompõe em
$n-2$ triângulos, e cada triângulo contribui com $180^\circ$,
a soma $S$ das amplitudes dos ângulos internos é
\[
\boxed{S=(n-2)\times180^\circ.}
\]

Na expressão, $n$ representa o número de lados,
$n-2$ o número de triângulos da decomposição,
e $180^\circ$ a soma das amplitudes dos ângulos
internos de cada triângulo. Por distributividade,
pode também escrever-se
\[
S=180^\circ n-360^\circ.
\]
As duas expressões são equivalentes.

\subsection*{2.4. Justificação para qualquer polígono convexo}

Considere-se um polígono convexo com $n$ lados,
sendo $n$ um número inteiro e $n\geq3$.
Escolha-se um dos seus vértices, por exemplo $A$.

Dos restantes $n-1$ vértices, dois são adjacentes a $A$:
os extremos dos dois lados que aí se encontram.
Os outros $n-3$ são não adjacentes, pelo que é possível
ligar $A$ a cada um deles através de uma diagonal.

Por o polígono ser \textbf{convexo}, essas diagonais
ficam no interior do polígono. Traçadas a partir do mesmo
vértice, não se cruzam no interior e dividem-no numa sequência
de regiões triangulares, sem sobreposições nem espaços por preencher.

Os $n-3$ segmentos interiores, juntamente com os dois lados
incidentes em $A$, originam $n-2$ setores triangulares:
\[
\underbrace{n-3}_{\text{diagonais}}
+\underbrace{1}_{\text{região triangular adicional}}
=\underbrace{n-2}_{\text{triângulos}}.
\]
Em particular, para $n=3$, não se traçam diagonais e
obtém-se o próprio triângulo, em conformidade com a expressão.

Como a soma das amplitudes dos ângulos internos dos triângulos
coincide com a soma das amplitudes dos ângulos internos do
polígono (questão 1.3), conclui-se que
\[
S=\underbrace{180^\circ+\cdots+180^\circ}_{n-2\text{ parcelas}}
=\boxed{(n-2)\times180^\circ}.
\]

\textbf{A justificação não depende de verificar alguns exemplos:}
explica por que razão a decomposição e a contagem dos triângulos
funcionam para qualquer polígono convexo, com $n\geq3$ lados.

\section*{Para concluir}

Uma forma de completar a frase da ficha é:

\begin{tcolorbox}[
colback=azulClaro,colframe=azulPDMI!40,boxrule=.5pt,
arc=2mm,left=8pt,right=8pt,top=6pt,bottom=6pt]
«Num polígono convexo com $n$ lados, a soma das amplitudes
dos ângulos internos é
\[
\boxed{S=(n-2)\times180^\circ}
\]
porque podemos dividir esse polígono em $n-2$ triângulos,
sem sobreposições nem espaços por preencher,
e a soma das amplitudes dos ângulos internos de cada
triângulo é $180^\circ$.»
\end{tcolorbox}

\section*{Se houver tempo -- polígono convexo com 12 lados}

Como $n=12$, a decomposição a partir de um vértice
origina $12-2=10$ triângulos. Logo,
\[
S=(12-2)\times180^\circ
=10\times180^\circ
=\boxed{1800^\circ}.
\]
Não é necessário desenhar o polígono, pois a justificação
da expressão geral garante a sua validade para $n=12$.

\endgroup


% ==========================================
% PARTE II - AULA ALTERNATIVA
% ==========================================

\partwithheader{Aula alternativa}

% PLANIFICACAO

\chapter{Planificação}

\begin{center}
{\large\bfseries\color{azulPDMI}Quanto precisa de virar o gato?}\\[3pt]
{\small Planificação da aula alternativa - Matemática, 7.º ano}
\end{center}

\section*{1. Objeto}
\textbf{Tema:} Geometria. \textbf{Tópico:} Soma das amplitudes dos ângulos internos de um polígono convexo.\\
\textbf{Objeto matemático central:} dedução e justificação da soma das amplitudes dos ângulos internos de um polígono convexo com $n$ lados.\\
\textbf{Proposta de duração:} 50 minutos.

\section*{2. Objetivos}

\subsection*{Objetivos gerais}

\subsection*{Objetivos específicos}

\section*{3. Pré-requisitos}

\section*{4. Recursos e materiais}

\section*{5. Desenvolvimento da aula}

\section*{6. Antecipação das dificuldades e respostas do professor}

\section*{7. Avaliação das aprendizagens}




% ENUNCIADO

\chapter{Enunciado}

\begin{center}
{\LARGE\bfseries\color{azul}
Quanto precisa de virar o gato?}\\[4pt]
{\normalsize Uma investigação sobre ângulos
internos e mudanças de direção}
\end{center}

\noindent
Nome: \rule{9cm}{.4pt}
\hfill
N.º: \rule{2cm}{.4pt}

\begin{desafio}
\textbf{O desafio}\\[3pt]

Um gato desloca-se em linha reta e,
sempre que chega a um vértice,
muda de direção para percorrer
o lado seguinte de um polígono.

\textbf{Como podemos descobrir quanto deve virar?}

E será possível usar essas viragens
para encontrar a soma das amplitudes
dos ângulos internos de qualquer
polígono convexo?

\end{desafio}

% =====================================
% PARTE I
% =====================================

\parte{Parte I}{Prever, experimentar e explicar}

Abre no \textbf{Scratch} o projeto
fornecido pelo professor.

O programa já está preparado:
não precisas de construir os blocos.

Sempre que premires a tecla
\textbf{Espaço}, o gato pergunta
\textbf{quantos lados} queres desenhar
e \textbf{quantos graus} deve virar
em cada vértice.

O gato efetua uma viragem depois
de percorrer \emph{cada} lado,
incluindo o último.

\begin{enumerate}[label=\textbf{1.\arabic*.}]

\item
Pretendes obter um quadrado.
Introduz \textbf{4 lados} e uma
viragem de \textbf{$80^\circ$}.

Antes de executar o desenho,
achas que o percurso vai fechar
um quadrado?
Explica a tua previsão.

\linhas{2}

\item
Testa a tua previsão.
Mantendo os quatro lados, procura
uma viragem que permita desenhar
um \textbf{quadrado}, sem alterar
o comprimento dos lados.

Regista o valor encontrado
e explica a tua escolha.

\linhas{2}

\item
No quadrado, a viragem e o
ângulo interno têm ambos
amplitude $90^\circ$.

Acontecerá o mesmo em todos
os polígonos?
Explica a tua previsão.

\linhas{2}

\end{enumerate}

\begin{nota}
\textbf{Para observar com atenção:}

Queremos desenhar o \emph{contorno de
um polígono regular convexo},
sem cruzamentos e percorrido
uma única vez.

Não basta que o gato termine
no ponto onde começou.
\end{nota}

% =====================================
% PARTE II
% =====================================

\parte{Parte II}{Investigar uma regularidade}

O gato mantém sempre o mesmo
comprimento de lado e a mesma
amplitude de viragem, em cada tentativa.

Investiga os polígonos
\textbf{regulares convexos} da tabela.

\begin{enumerate}[label=\textbf{2.\arabic*.}]

\item
Para cada número de lados, descobre
a viragem que faz o gato desenhar
o polígono pretendido.

Regista o valor e calcula a
\textbf{soma de todas as viragens}
realizadas numa volta completa.

\end{enumerate}

\renewcommand{\arraystretch}{1.6}

\begin{tabularx}{\linewidth}{
>{\raggedright\arraybackslash}X
>{\centering\arraybackslash}p{1.65cm}
>{\centering\arraybackslash}p{2.7cm}
>{\centering\arraybackslash}p{2.8cm}}

\rowcolor{azulclaro}

\textbf{Polígono} &
\textbf{Lados} &
\textbf{Viragem por vértice} &
\textbf{Soma das viragens}
\\ \hline

Triângulo equilátero & 3 & & \\
Quadrado & 4 & & \\
Pentágono regular & 5 & & \\
Hexágono regular & 6 & & \\

\end{tabularx}

\renewcommand{\arraystretch}{1.55}

\begin{enumerate}[
label=\textbf{2.\arabic*.},
start=2]

\item
O que muda quando aumenta
o número de lados?
O que se mantém?

Compara os resultados com
os de um colega.

\linhas{2}

\item
Se o gato desenhasse um polígono
\textbf{regular convexo} com $n$ lados
e efetuasse uma viragem de amplitude
$v$ em cada vértice, que relação
esperarias encontrar entre $n$ e $v$?

Explica.

\linhas{2}

\end{enumerate}

% =====================================
% PARTE III
% =====================================

\parte{Parte III}{}

Prime agora a tecla \textbf{I}.

O gato desenha um
\textbf{paralelogramo que não é
um retângulo}.

Observa o percurso e os valores
das viragens nos blocos do programa.

\begin{enumerate}[label=\textbf{3.\arabic*.}]

\item
Indica as quatro viragens
e calcula a sua soma.

\linhas{1}

\item
A soma obtida é igual à dos
polígonos regulares?

O que mostra este exemplo?

Será suficiente para teres
a certeza de que acontece
o mesmo em \emph{todos}
os polígonos convexos?

\linhas{2}

\end{enumerate}

% =====================================
% PARTE IV
% =====================================

\Needspace{12\baselineskip}

\parte{Parte IV}{}

A viragem realizada pelo gato
num vértice corresponde à amplitude
de um \textbf{ângulo externo}
do polígono.

Repara na relação entre esse ângulo
e o ângulo interno do mesmo vértice.

\begin{center}

\begin{tikzpicture}[scale=.95,>=Stealth]

\coordinate (B) at (0,0);
\coordinate (C) at (2.2,1.6);

\fill[verde!13]
(B)--(1.05,0)
arc (0:36:1.05)--cycle;

\fill[azul!13]
(B)--({.98*cos(36)},{.98*sin(36)})
arc (36:180:.98)--cycle;

\draw[very thick,azul,
-{Stealth[length=2.8mm]}]
(-3,0)--(-.12,0);

\draw[very thick,azul,
-{Stealth[length=2.8mm]}]
(B)--(C);

\draw[gray,dashed,thick]
(B)--(2.7,0);

\draw[verde!75!black]
(1.05,0) arc (0:36:1.05);

\draw[azul!90!black]
({.98*cos(36)},{.98*sin(36)})
arc (36:180:.98);

\node at (-2.85,-.22) {A};
\node[below] at (B) {B};
\node[above] at (2.2,1.6) {C};

\node[verde!60!black]
at (1.44,.26) {$\beta$};

\node[azul!90!black]
at (-.48,.72) {$\alpha$};

\node[gray,below]
at (2.12,-.10) {prolongamento};

\end{tikzpicture}

\end{center}

\begin{enumerate}[label=\textbf{4.\arabic*.}]

\item
Que relação existe entre a
amplitude $\alpha$ do ângulo interno
e a amplitude $\beta$ do ângulo
externo correspondente?

Justifica.

\linhas{1}

\item
Num pentágono convexo existem
cinco pares de ângulos,
um interno e um externo adjacente,
em cada vértice.

Quanto somam, ao todo,
os \textbf{cinco pares}?

Como poderás encontrar a soma
dos cinco ângulos internos?

\linhas{2}

\item
Generaliza o teu raciocínio
para um polígono convexo com
$n$ lados e propõe uma expressão
para a soma $S$ dos seus
ângulos internos.

Explica o significado
de cada parcela.

\linhas{2}
\newpage
\item
Os exemplos do Scratch sugerem
que a soma das viragens é $360^\circ$.

\textbf{Por que razão isso acontece
ao percorrer, uma única vez e sempre
no mesmo sentido, o contorno de
qualquer polígono convexo?}

Explica sem te baseares
apenas nos exemplos.

\linhas{2}

\end{enumerate}

% =====================================
% CONCLUSAO
% =====================================

\begin{conclusao}

\textbf{A nossa conclusão}

Para um polígono convexo com $n$ lados,
a soma das amplitudes dos seus ângulos
internos é

\[
\boxed{S=\rule{5.2cm}{.35pt}}
\]

A expressão é válida porque
\rule{8.8cm}{.35pt}

\vspace{5pt}

\rule{\linewidth}{.35pt}

\end{conclusao}

\vspace{1mm}

{\footnotesize\itshape
Desafio facultativo:
no programa 1, experimenta 5 lados
e uma viragem de $144^\circ$.

O gato regressa ao início,
mas desenha um pentágono convexo?

O que te ensina esta experiência?
}


% RESOLUCAO

\chapter{Resolução}

\begin{center}
{\large\bfseries\color{azulPDMI}Quanto precisa de virar o gato?}\\[2pt]
{\small Resolução das questões 1.1 a 4.4, conclusão e desafio facultativo}
\end{center}

\begin{nota}
\textbf{Condições consideradas.} Estuda-se o contorno de polígonos
\textbf{convexos}, percorrido uma única vez, sem cruzamentos e sempre
no mesmo sentido. Em cada vértice considera-se \textbf{um} ângulo
externo, correspondente à mudança de direção do gato. O programa
faz uma viragem depois de cada lado, \textbf{incluindo o último}.
Mover o gato com a caneta levantada para o centro do palco não integra
o percurso geométrico.
\end{nota}

\section*{Parte I}

\noindent\textbf{1.1. Quatro lados e uma viragem de $80^\circ$}

Não se obtém um quadrado. Num quadrado, os lados consecutivos são
perpendiculares, pelo que a mudança de direção em cada vértice tem
de ser de $90^\circ$. Se o gato virar apenas $80^\circ$, a direção
do lado seguinte não será perpendicular à do lado anterior. A
sequência de quatro segmentos de igual comprimento não fecha o
contorno de um quadrado.

A soma das viragens programadas é
\[
4\times80^\circ=320^\circ,
\]
isto é, $40^\circ$ menos do que uma volta completa. Esta conta ajuda
a interpretar o resultado, mas a razão por que o contorno de um
polígono convexo corresponde a uma volta completa será justificada
na questão~4.4; não deve ser simplesmente pressuposta no lançamento
da tarefa.

\begin{center}
\begin{minipage}{.46\linewidth}
\centering
\begin{tikzpicture}[scale=.72,>=Stealth]
  \draw[thick,azulPDMI,line join=round]
  (0,0)--(2,0)--(2.347,-1.970)--(.468,-2.654)--(-.532,-.922);
  \fill[azulPDMI] (0,0) circle (2.5pt);
  \fill[orange!85!black] (-.532,-.922) circle (2.5pt);
  \node[above left] at (0,0) {$A$};
  \node[left] at (-.532,-.922) {$A'$};
  \draw[gray,dashed] (-.532,-.922)--(0,0);
\end{tikzpicture}\\[-3pt]
{\footnotesize $4$ viragens de $80^\circ$: $A'\neq A$.}
\end{minipage}\hfill
\begin{minipage}{.46\linewidth}
\centering
\begin{tikzpicture}[scale=.72,>=Stealth]
  \draw[thick,azulPDMI,line join=round]
    (0,0)--(2,0)--(2,-2)--(0,-2)--cycle;
  \fill[azulPDMI] (0,0) circle (2.5pt);
  \node[above left] at (0,0) {$A$};
\end{tikzpicture}\\[-3pt]
{\footnotesize $4$ viragens de $90^\circ$: quadrado.}
\end{minipage}
\end{center}

\noindent\textbf{1.2. Viragem necessária para desenhar o quadrado}

A amplitude correta é $\boxed{90^\circ}$. Como o gato anda sempre
o mesmo número de passos, os quatro lados têm o mesmo comprimento.
Ao realizar uma viragem de $90^\circ$ depois de cada lado, obtêm-se
quatro ângulos retos e o contorno fecha sem cruzamentos. A soma das
quatro mudanças de direção é
\[
90^\circ+90^\circ+90^\circ+90^\circ=360^\circ.
\]

\noindent\textbf{1.3. Comparação entre ângulo interno e viragem}

A coincidência verificada no quadrado \textbf{não é geral}. Num
triângulo equilátero, por exemplo, cada ângulo interno mede $60^\circ$,
mas o gato tem de virar $120^\circ$ ao passar de um lado ao seguinte.
Os dois ângulos são adjacentes e suplementares:
\[
60^\circ+120^\circ=180^\circ.
\]
Num pentágono regular, o ângulo interno mede $108^\circ$ e a
viragem mede $72^\circ$. A igualdade entre a viragem e o ângulo
interno no quadrado é, portanto, um caso particular.

\begin{nota}
\textbf{Resposta possível de um aluno.} «Não. No quadrado os dois
ângulos medem $90^\circ$, mas num triângulo equilátero o gato vira
$120^\circ$ e o ângulo interno mede $60^\circ$. Juntos perfazem
$180^\circ$.»
\end{nota}

\section*{Parte II}

\noindent\textbf{2.1. Tabela completa}

\begin{center}
\renewcommand{\arraystretch}{1.32}
\begin{tabularx}{\linewidth}{|>{\raggedright\arraybackslash}X|
>{\centering\arraybackslash}p{1.5cm}|
>{\centering\arraybackslash}p{2.6cm}|
>{\centering\arraybackslash}p{2.65cm}|}
\hline
\rowcolor{azulClaro}
\textbf{Polígono} & \textbf{N.º de lados} &
\textbf{Viragem por vértice} & \textbf{Soma das viragens}\\\hline
Triângulo equilátero & $3$ & $120^\circ$ & $3\times120^\circ=360^\circ$\\\hline
Quadrado & $4$ & $90^\circ$ & $4\times90^\circ=360^\circ$\\\hline
Pentágono regular & $5$ & $72^\circ$ & $5\times72^\circ=360^\circ$\\\hline
Hexágono regular & $6$ & $60^\circ$ & $6\times60^\circ=360^\circ$\\\hline
\end{tabularx}
\end{center}

Em cada caso, o valor indicado é a \textbf{mudança de direção},
não a amplitude de um ângulo interno. Os valores correspondem ao
percurso do contorno do polígono regular convexo, sem cruzamentos e
com uma viragem em cada vértice. A mera igualdade
$n\,v=360^\circ$ não garante, sem outras condições, que um
percurso arbitrário seja o contorno simples pretendido.

\noindent\textbf{2.2. O que muda e o que se mantém?}

À medida que o número de lados aumenta de $3$ para $6$, aumenta o
número de viragens, mas diminui a amplitude de cada uma: $120^\circ$,
$90^\circ$, $72^\circ$ e $60^\circ$. A \textbf{soma} das amplitudes das
viragens mantém-se igual a $360^\circ$. Nos polígonos regulares,
todas as viragens de um mesmo percurso têm a mesma amplitude.

É importante distinguir três quantidades: \emph{número de lados},
\emph{amplitude de cada viragem} e \emph{soma das viragens}. Esta
última é a quantidade que se mantém nos casos estudados.

\noindent\textbf{2.3. Generalização para $n$ lados}

Se o gato percorre uma única vez o contorno de um polígono regular
convexo com $n$ lados, efetuando uma viragem de amplitude $v$ em cada
vértice, efetua $n$ viragens iguais. Como perfaz uma volta completa,
\[
\underbrace{v+v+\cdots+v}_{n\text{ parcelas}}=360^\circ
\quad\Longleftrightarrow\quad
n\,v=360^\circ.
\]
Logo,
\[
\boxed{v=\frac{360^\circ}{n}},\qquad n\in\mathbb{N},\ n\geq3.
\]
Esta relação dá a amplitude de \textbf{cada} ângulo externo de um
\textbf{polígono regular convexo}. Não significa que todos os ângulos
externos de um polígono convexo qualquer sejam iguais.

\section*{Parte III}

\noindent\textbf{3.1. Viragens no paralelogramo}

Ao premir a tecla \textbf{I}, o gato desenha um paralelogramo não
retângulo com quatro mudanças de direção, por esta ordem:
\[
60^\circ,\quad120^\circ,\quad60^\circ,\quad120^\circ.
\]
A sua soma é
\[
60^\circ+120^\circ+60^\circ+120^\circ=\boxed{360^\circ}.
\]
As amplitudes dos ângulos internos correspondentes são,
respetivamente, $120^\circ$, $60^\circ$, $120^\circ$ e $60^\circ$.

\noindent\textbf{3.2. Comparação e alcance da experiência}

Sim, a soma é igual à obtida para os polígonos regulares: $360^\circ$.
O paralelogramo mostra que \textbf{a igualdade das amplitudes das
viragens não é necessária} para se obter essa soma: há duas viragens
de $60^\circ$ e duas de $120^\circ$.

Contudo, \textbf{um paralelogramo não demonstra} que a propriedade
é válida para todos os polígonos convexos. A experiência fornece
um novo exemplo favorável à conjetura, mas é necessário apresentar
um argumento matemático que abranja qualquer polígono convexo.

\begin{nota}
\textbf{Resposta possível de um aluno.} «A soma continua a ser
$360^\circ$, apesar de o gato não virar sempre o mesmo. Mas testar
este paralelogramo não chega para saber o que acontece com todos os
polígonos convexos.»
\end{nota}

\section*{Parte IV}

\noindent\textbf{4.1. Relação entre o ângulo interno e o externo}

O ângulo externo $\beta$ considerado na tarefa é formado pelo
prolongamento de um lado do polígono e pelo lado seguinte. O ângulo
interno $\alpha$ e esse ângulo externo têm um lado comum; os
outros lados são semirretas opostas. São, portanto, \textbf{ângulos
adjacentes suplementares}, pelo que
\[
\boxed{\alpha+\beta=180^\circ}
\qquad\text{ou}\qquad
\boxed{\alpha=180^\circ-\beta}.
\]
A mudança de direção do gato corresponde ao ângulo externo,
\textbf{não}, em geral, ao ângulo interno do polígono.

\noindent\textbf{4.2. O pentágono convexo}

Um pentágono tem cinco vértices e, por isso, cinco pares de um
ângulo interno com o respetivo externo adjacente. Cada par soma
$180^\circ$; assim, a soma das amplitudes dos \textbf{dez ângulos}
(cinco internos e cinco externos) é
\[
5\times180^\circ=900^\circ.
\]
A soma dos cinco ângulos externos é $360^\circ$. Logo, a soma dos
cinco ângulos internos é
\[
\boxed{S_5=900^\circ-360^\circ=540^\circ}.
\]
A resposta não depende de o pentágono ser regular. Se fosse regular,
teríamos cinco ângulos internos iguais a $108^\circ$ e verificar-se-ia
$5\times108^\circ=540^\circ$.

\noindent\textbf{4.3. Generalização para um polígono com $n$ lados}

Considere-se um polígono convexo com $n\geq3$ lados. Designemos por
$\alpha_i$ a amplitude do ângulo interno do vértice $i$ e por
$\beta_i$ a amplitude do ângulo externo correspondente, escolhido
sempre no mesmo sentido do percurso. Como cada par é suplementar,
\[
\alpha_i+\beta_i=180^\circ
\qquad(i=1,\ldots,n).
\]
Somando estas $n$ igualdades, obtém-se
\[
\underbrace{\sum_{i=1}^{n}\alpha_i}_{S}
+\underbrace{\sum_{i=1}^{n}\beta_i}_{360^\circ}
=n\times180^\circ.
\]
Consequentemente,
\[
\boxed{S=n\times180^\circ-360^\circ=(n-2)\times180^\circ}.
\]
Na expressão, $n\times180^\circ$ é a soma das amplitudes dos
$n$ pares de ângulos suplementares; subtraem-se $360^\circ$ porque
essa é a soma das amplitudes dos ângulos externos. O resultado é a
soma das amplitudes dos ângulos \textbf{internos} do polígono.

Esta dedução não exige que os ângulos internos sejam todos iguais;
por isso, aplica-se aos polígonos convexos regulares e não regulares.

\noindent\textbf{4.4. Justificação geométrica da soma dos ângulos externos}

A explicação não deve limitar-se à afirmação de que o gato
«regressou ao sítio de onde partiu». Na verdade, a orientação final
poderia coincidir com a inicial após zero, uma ou várias voltas
completas; e um desenho poderia fechar com autointerseções.

A propriedade resulta das condições geométricas do percurso:
\begin{enumerate}[label=\textbf{(\roman*)},leftmargin=1.15cm,itemsep=2pt,topsep=3pt]
\item o gato percorre \textbf{uma única vez} o contorno de um polígono
\textbf{convexo}, sem o atravessar;
\item em cada vértice muda de direção \textbf{sempre no mesmo sentido},
pois todos os ângulos externos considerados são viragens exteriores;
\item ao acompanhar, por ordem, as orientações dos lados de um
contorno convexo, a direção percorre as orientações de uma
\textbf{única volta completa}, sem retroceder nem fazer uma segunda volta.
\end{enumerate}

Assim, a soma das mudanças de direção, incluindo a última viragem
que recupera a orientação inicial, é
\[
\boxed{\beta_1+\beta_2+\cdots+\beta_n=360^\circ}.
\]

\textbf{Uma forma de tornar o argumento visível na discussão:}
imaginar os vetores orientados de cada lado transportados, sem
rodar, para um mesmo ponto. Como o contorno é convexo, as direções
sucedem-se pela mesma ordem e os ângulos entre direções consecutivas
preenchem exatamente uma volta em torno desse ponto. Não há
sobreposições de voltas nem direções percorridas em sentido inverso.
Este argumento geométrico explica a regularidade para \textbf{qualquer}
polígono convexo, em vez de depender de uma lista finita de testes.

\begin{nota}
\textbf{Precisão matemática.} A afirmação de que as direções dos
lados perfazem exatamente uma volta usa a \emph{convexidade} e o
facto de o contorno ser percorrido uma vez. Não decorre apenas de o
gato terminar na mesma posição ou orientação. Uma justificação
mais formal pode recorrer à variação total de direção num contorno
convexo; ao nível do 7.º ano, o argumento por uma volta completa,
apoiado na representação das direções, é a explicação geométrica
adequada a explorar e discutir.
\end{nota}

\section*{Conclusão a registar na ficha}

\begin{conclusao}
Para um polígono convexo com $n$ lados, a soma das amplitudes dos
seus ângulos internos é
\[
\boxed{S=(n-2)\times180^\circ}.
\]
\textbf{Justificação:} em cada vértice, o ângulo interno e o externo
adjacente somam $180^\circ$; existem $n$ pares desse tipo. Ao
percorrer uma vez o contorno de um polígono convexo, a soma dos
ângulos externos é $360^\circ$. Por isso, a soma dos ângulos
internos obtém-se subtraindo $360^\circ$ a $n\times180^\circ$.
\end{conclusao}

\Needspace{14\baselineskip}
\section*{Desafio facultativo}

O gato desenha uma \textbf{estrela de cinco pontas} (pentagrama),
em que os segmentos se cruzam. Ao fim dos cinco segmentos, volta
à posição inicial, mas o desenho \textbf{não é o contorno simples
de um pentágono convexo}. A soma das cinco viragens programadas é
\[
5\times144^\circ=\boxed{720^\circ}=2\times360^\circ.
\]
O gato completou \textbf{duas voltas na orientação}, embora tenha
regressado ao ponto de partida. Logo, fechar um percurso não é
condição suficiente para garantir que a soma das viragens seja
$360^\circ$. É indispensável explicitar as condições de
\emph{convexidade}, \emph{contorno simples} e \emph{uma única volta}.

\begin{center}
\begin{tikzpicture}[scale=.87,line join=round]
\foreach \k in {0,...,4}{
  \coordinate (V\k) at ({90-72*\k}:1.85);
}
\draw[very thick,azulPDMI]
 (V0)--(V2)--(V4)--(V1)--(V3)--cycle;
\fill[azulPDMI] (V0) circle (2.2pt);
\node[above] at (V0) {Início / fim};
\end{tikzpicture}
\end{center}

\section*{Síntese da resolução e cuidados na discussão}

\begin{enumerate}[label=\textbf{\arabic*.},leftmargin=1.05cm,itemsep=3pt,topsep=3pt]
\item \textbf{O que o Scratch permite observar:} valores de viragem
para polígonos regulares e um exemplo não regular em que a soma
exterior continua a ser $360^\circ$.
\item \textbf{O que precisa de ser justificado:} por que razão a
mudança total de direção, num único percurso do contorno convexo,
é uma volta completa.
\item \textbf{A passagem essencial:} a relação
$\alpha_i+\beta_i=180^\circ$ em cada um dos $n$ vértices, seguida
da subtração da soma dos ângulos externos.
\item \textbf{Um erro a prevenir:} escrever $v=360^\circ/n$
para \emph{cada} ângulo externo de um polígono não regular. Essa
igualdade só é válida quando as viragens são todas iguais.
\item \textbf{Outro erro a prevenir:} concluir que a soma é $360^\circ$
por o gato ter regressado ao ponto inicial; a estrela mostra por
que esse argumento é insuficiente.
\item \textbf{Articulação com a primeira aula:} chega-se à mesma
expressão $S=(n-2)\times180^\circ$, mas por ângulos externos e
suplementaridade, \textbf{sem decompor o polígono em triângulos}.
\end{enumerate}


% ==========================================
% FIM DO DOCUMENTO
% ==========================================

\end{document}

